\section{Discretisation}

We consider a mesh $\mathcal{M}$ over a surface $\Omega$ with polygons; along with the $\bm{x}$ spacial variable. Let $j$ be a polygonal cell of this mesh and $\ell$ an edge $[r, s]$. We will denote $\bm{n}_{\ell j}$ the outward from cell $j$ normal at the edge $\ell$. The centroid of the cell $j$ is $\bm{x}_j$, the middle of the edge $\ell$ is $\bm{x}_\ell$ and the position of the node of $r$ (resp. s) is called $\bm{x}_r$ (resp. $\bm{x}_s$). A second cell $k$ is said to be a neighbouring cell of $j$, if $\exists ! \ell \in \mathcal{M} : (\ell \in k) \cap (\ell \in j)$. The boudnaries of the 

\begin{figure}[H]
  \centering
	\begin{tikzpicture}[scale=1.5, solidnode/.style={circle,fill=black,inner sep=1.6pt},opencircle/.style={circle,draw=black,very thin,inner sep=1.6pt}, bar/.style={line width=0.7pt,draw=black},arrowstyle/.style={-Stealth,thin}]

 \coordinate (a) at (0.5,0.5);
 \coordinate (b) at (0.7,1.7);
 \coordinate (c) at (1.8,2);
 \coordinate (d) at (2.8,1.5);
 \coordinate (e) at (2.1,0.4);

 \coordinate (ell) at ($(a)!0.5!(b)$);

    \draw[thick] (a) -- (b) -- (c) -- (d) -- (e) -- cycle;
    \fill[color=cyan!40!gray] (a) -- (b) -- (c) -- (d) -- (e) -- cycle;
\foreach \p in {a,b,c,d, e}
			\node[opencircle, fill=white] at (\p) {};

	\draw (a) node [left] {$r$};
	\draw (b) node [left] {$s$};

	\draw (ell) node [below left] {$\ell$};
	\draw (ell) node [] {$\times$};
  \coordinate (nell) at ($ (ell)!0.5!90:(b) $);
  \draw [arrowstyle](ell)--(nell) node [above right,inner sep=1pt] {\small $\bm{n}_{\ell}$};
\end{tikzpicture}
\caption{A cell $j$, an edge $\ell$ and its normal $\bm{n}_\ell$}
  \label{fig:example_cell}
\end{figure}
We start by recalling divergence theorem. Given a field $\bm{u}$ and a closed surface $V$,
\be
\int_V \nabla \cdot \bm{u}\, dV = \int_{\partial V} \bm{n} \cdot \bm{u} \, dS
\label{divergence_theorem}
\ee
where $\partial V$ denotes the boundaries of the closed surface $V$ and $\bm{n}$ the outpointing normal.

We define the field $\bm{u}^\star$, so that
\be
\dis \frac{\partial \bm{u}^\star}{\partial t} + (\bm{u}^\star \cdot \nabla) \bm{u} - \nu \nabla^2 \bm{u}^\star - \bm{g} = 0.
\label{define_ustar}
\ee
which is, in essence, the same equation as \eqref{navier_stokes} but with the pressure field term omitted after applying a semi-implicit (resp. implicit) discretization of the and transport (resp. diffusion) term.

We will then integrate \eqref{define_ustar} over a cell $j$ (see figure \ref{fig:example_cell}) and discretize in time.
\be
\dis \int_j \frac{\bm{u}^\star - \bm{u}}{\Delta t} + (\bm{u}^\star \cdot \nabla) \bm{u} - \nu \nabla^2 \bm{u}^\star - \bm{g} = 0.
\label{integrate_ustar}

\ee
We now split this last equation and evalute term by term.

\subsection{Constant term}

\be
\int_j \bm{g} = V_j \bm{g}
\ee

\subsection{Convection term}
The convection term can be expressed using the divergence theorem \eqref{divergence_theorem},
\be
\dis \int_j (\bm{u}^\star \cdot \nabla) \bm{u} 
= \int_{\partial j} (\bm{u}^\star \cdot \bm{n}_{\partial j}) \bm{u}
= \sum_{\ell \in j} \int_\ell (\bm{u}^\star \cdot \bm{n}_{\ell j}) \bm{u}
= \sum_{\ell \in j} |\ell| (\bm{u}_\ell^\star \cdot \bm{n}_{\ell j}) \bm{u}_\ell.
\ee
Please 

% - \nu \nabla^2 \bm{u}^\star - \bm{g} = 0.
